\section*{Chapitre \ref{ch:ll:rust-for-low-latency} --- Rust pour la faible latence}

\begin{solution}{exo:ll:rust-for-low-latency:1}
Sur la mesure retenue, environ 10.3 et \qty{10.6}{\nano\second}~: Rust reste à quelques pour cent de C++ (un cinquième au plus
d'un passage à l'autre). À 5 millions de messages par seconde, un écart de \qty{0.3}{\nano\second} représente 1.5 millisecondes par
seconde, un bruit négligeable face aux autres coûts du chemin critique.
\end{solution}

\begin{solution}{exo:ll:rust-for-low-latency:2}
\texttt{\textquotedbl{}ESZ6\textquotedbl{}.to\_string()} alloue (une \texttt{String} ne dispose pas de tampon interne)~; \texttt{Vec::new()} n'alloue pas~; son premier
\texttt{push} alloue (place pour quatre petits éléments)~; \texttt{[0u8; 24]} se trouve sur la pile~; \texttt{Box::new(3)} alloue.
\end{solution}

\begin{solution}{exo:ll:rust-for-low-latency:3}
Rust~: la \texttt{String} du symbole (1, là où le tampon interne de C++ l'évitait), la \texttt{String} de l'identifiant (1, comme en
C++), le vecteur (1~: le premier push réserve quatre places, contre trois étapes de croissance en C++), l'entrée de la table
(0~: la clé est déplacée à l'intérieur et la table était déjà allouée, contre un nœud et une copie de clé en C++), la copie des
exécutions (1, comme la copie de la lambda en C++), la fermeture encadrée (1, contre deux dans \texttt{std::function}). $1+1+1+0+1+1 = 5$ contre $0+1+3+2+1+2 = 9$.
\end{solution}

\begin{solution}{exo:ll:rust-for-low-latency:4}
Quatre~: deux chargements de 64 bits de deux \texttt{u32} chacun, élargis à 64 bits et additionnés dans deux accumulators vectoriels à deux voies.
Le reste de zéro à trois éléments passe par la boucle scalaire en \texttt{.LBB16\_4}.
\end{solution}

\begin{solution}{exo:ll:rust-for-low-latency:5}
Chaque itération attend le chargement d'une adresse dispersée~; la vérification est une comparaison indépendante avec un branchement prédit,
exécutée dans l'ombre de cette attente. Elle se verrait dans une boucle limitée par le processeur plutôt que par la mémoire, sur des données présentes dans le
cache de premier niveau, et surtout là où la vérification empêche la vectorisation.
\end{solution}

\begin{solution}{exo:ll:rust-for-low-latency:6}
Non. \texttt{CheckedIndex::new(\&[99], 100)} suivi de \texttt{sum(\&v, 100)} avec \texttt{v} de longueur 10 lirait
\texttt{v[99]} hors des bornes, en code sûr~: l'abstraction aurait fait confiance à son appelant pour \texttt{len}. La comparaison
avec la longueur réelle de la tranche est ce qui la rend \omterm{def:ll:rust-for-low-latency:sound}{saine}.
\end{solution}

\begin{solution}{exo:ll:rust-for-low-latency:7}
Le \texttt{match} explicite rend visible le cas d'erreur~: \texttt{flatten()} ignore silencieusement une erreur de décodage, alors que le
\texttt{match} peut la compter et s'arrêter. Aucune des deux versions n'alloue (le décodeur emprunte le tampon), et aucune ne contient
de chemin de panique dans le corps de la boucle lorsque les longueurs sont vérifiées par les comparaisons propres au décodeur~; la
différence porte sur le comportement face aux entrées erronées, non sur la vitesse.
\end{solution}

\begin{solution}{exo:ll:rust-for-low-latency:8}
Un gestionnaire de données de marché est un flux unique et ordonné~: il ne peut pas être réparti sur plusieurs cœurs sans perdre l'ordre, et un ordonnanceur à
vol de tâches ajoute des files d'attente, des migrations entre threads et des \omterm{def:ll:memory-hierarchy-and-caches:miss}{défauts de cache}, exactement la \omterm{def:ll:where-latency-comes-from:tail}{gigue} que le chemin critique évite. Utiliser un thread épinglé
qui interroge en boucle son socket ou son anneau~; réserver async aux nombreuses connexions lentes.
\end{solution}

\begin{solution}{pb:ll:rust-for-low-latency:1}
\begin{enumerate}
\item Environ \qty{10.3}{\nano\second} (C++) et \qty{10.6}{\nano\second} (Rust) sur la mesure retenue~: un rapport d'environ
1.03, et un écart d'un cinquième au plus d'un passage à l'autre.
\item LLVM vectorise la boucle à compteur à \texttt{-O}~; g++ 11 ne vectorise pas à \texttt{-O2} (chapitre~14).
\item \texttt{-O3}, ou g++ 12 ou une version ultérieure à \texttt{-O2}, ou une boucle vectorisée explicitement.
\item Qu'elles ne coûtent rien de mesurable à cet endroit~: les collectes vérifiées et non vérifiées prennent le même temps,
à l'intérieur du bruit, dans les deux langages.
\item Cinq (Rust) et neuf (C++) pour les ébauches~; aucune pour l'un ou l'autre gestionnaire corrigé.
\item \texttt{String}, qui alloue toujours, même pour un symbole court.
\item Un \omterm{def:ll:rust-for-low-latency:alloc}{allocateur global} de comptage installé dans le binaire de test, et une assertion sur le nombre d'allocations par message après la phase de chauffe.
\item Le compteur monte et le test échoue, à condition que le test sollicite les chemins critiques de la dépendance avec des données de forme réaliste.
\item Les \omterm{def:ll:rust-for-low-latency:race}{courses aux données}, les utilisations après libération et les références pendantes, ainsi que les lectures hors bornes (transformées en paniques).
\item Dans les appels à la bibliothèque C (affinité), dans les structures sans verrou (chapitre~11) et dans l'indexation non vérifiée
lorsqu'une \omterm{def:ll:rust-for-low-latency:sound}{abstraction saine} la justifie.
\item Qu'aucune suite d'appels passant par son interface sûre ne puisse provoquer de \omterm{def:ll:cpp-for-latency-iii-the-compiler:ub}{comportement indéfini}.
\item Un bogue a été détecté avant de corrompre la mémoire~; le processus doit cesser de coter (annuler les ordres, chapitre~22) et
redémarrer, et non poursuivre.
\item \textbf{Résultat nommé.} Le décodeur Rust consomme environ 1.03 fois le temps par message du décodeur C++ sur la mesure retenue
(un écart d'un cinquième au plus d'un passage à l'autre)~; les ébauches allouent cinq fois (Rust) et neuf fois (C++) par message, les gestionnaires
corrigés n'allouent jamais.
\item Non~: l'écart reste dans le bruit d'un passage à l'autre et demeure petit face aux autres étapes du chemin critique.
\item Le décodeur de flux et le constructeur de carnet~: analyse avec état d'une entrée non fiable, là où les erreurs mémoire sont les plus probables et les plus
coûteuses.
\item Autour du chemin critique~: connexions REST et WebSocket vers les plateformes, flux drop copy, services internes.
\item Par le fichier de plan écrit par \texttt{firm.coreplan} et lu par les deux jumeaux \texttt{firm.affinity}, et par des
jeux de référence que les deux implémentations doivent reproduire.
\item Le code C++ existant, l'outillage et l'expérience acquise~; certaines bibliothèques~; et le temps consacré au portage au détriment des fonctionnalités.
\item Exécuter les deux sur les mêmes entrées enregistrées et comparer les sorties message par message (tests différentiels, chapitre~25).
\item Elle déplace les erreurs de mémoire et de concurrence de la production vers le compilateur, à la même vitesse.
\end{enumerate}
\end{solution}

\begin{solution}{iq:ll:rust-for-low-latency:1}
L'abstraction se compile en ce qu'on écrirait à la main. La vérifier en lisant l'assembleur optimisé de la chaîne et de la
boucle écrite à la main, et en mesurant les deux sur des données réalistes.
\iqlookfor{assembleur et mesure, non la foi.}
\end{solution}

\begin{solution}{iq:ll:rust-for-low-latency:2}
Rarement~: une vérification prédite coûte peu, et le compilateur supprime celles qu'il peut prouver superflues. Elles coûtent
lorsqu'elles empêchent la vectorisation ou se trouvent dans une boucle limitée par le processeur. Les supprimer sans risque avec
des itérateurs, avec une assertion sur la longueur avant la boucle, ou derrière une \omterm{def:ll:rust-for-low-latency:sound}{abstraction saine} qui valide les indices une seule fois.
\iqlookfor{élimination par preuve et accès non vérifié sûr.}
\end{solution}

\begin{solution}{iq:ll:rust-for-low-latency:3}
Son interface sûre ne peut pas servir à provoquer un \omterm{def:ll:cpp-for-latency-iii-the-compiler:ub}{comportement indéfini}~: chaque invariant dont le code unsafe a besoin est établi et
vérifié par l'abstraction elle-même, et non supposé chez les appelants~; chaque bloc unsafe énonce son argument.
\iqlookfor{invariants pris en charge par l'abstraction.}
\end{solution}

\begin{solution}{iq:ll:rust-for-low-latency:4}
Les références mutables sont exclusives, et seuls les types \texttt{Send}/\texttt{Sync} traversent les threads, ce que le compilateur vérifie à la compilation. Il n'empêche pas
les conditions de course dans la logique (deux opérations atomiques qui auraient dû n'en faire qu'une), les interblocages, ni les courses
à l'intérieur du code \texttt{unsafe}.
\iqlookfor{\omterm{def:ll:rust-for-low-latency:race}{courses aux données} et conditions de course.}
\end{solution}

\begin{solution}{iq:ll:rust-for-low-latency:5}
Non~: un moteur d'exécution ordonnance les tâches sur plusieurs threads et les entremêle, ce qui ajoute de la latence et de la \omterm{def:ll:where-latency-comes-from:tail}{gigue}~;
le chemin critique est un thread épinglé qui interroge en boucle son entrée. Async convient aux nombreuses connexions en attente autour de lui.
\iqlookfor{\omterm{def:ll:where-latency-comes-from:tail}{gigue} d'ordonnancement contre thread dédié.}
\end{solution}

\begin{solution}{iq:ll:rust-for-low-latency:6}
Installer un \omterm{def:ll:rust-for-low-latency:alloc}{allocateur global} de comptage dans un binaire de test (il voit toutes les allocations, y compris celles des
dépendances), solliciter le composant avec des entrées de forme réaliste au-delà de la phase de chauffe, et affirmer que le compteur
ne change pas~; exécuter le test en intégration continue.
\iqlookfor{l'\omterm{def:ll:rust-for-low-latency:alloc}{allocateur global} comme point d'étranglement unique.}
\end{solution}
